%% title: Hausman Specification Testing and the LM Test by Artificial Regression
\documentclass[11pt]{article}
\usepackage{metricsstyle}

\begin{document}

\handoutheader{MLE V}{Supplementary lecture on maximum likelihood: Hausman specification testing (different from W, LM, LR), and the LM test computed by an artificial regression, with the test for AR(1) errors as the example}{}

\runin{Previously} \emph{MLE~IV}: restricted ML $\tilde\theta$ (under $H_{0}\colon r(\theta)=0$, $r$
restrictions, $R(\theta)=\partial r/\partial\theta'$), unrestricted ML $\hat\theta$, the information matrix
$\mathcal I$, and
$\operatorname{var}\bigl(\sqrt n(\tilde\theta-\theta_{0})\bigr)=\mathcal I^{-1}-\mathcal I^{-1}R'(R\,\mathcal I^{-1}R')^{-1}R\,\mathcal I^{-1}$.
Throughout, ``$\operatorname{var}$'' of an estimator means its asymptotic variance.

\section{Hausman specification testing}

Consider two estimators of $\theta$:
\begin{flatlist}
 \item $\tilde\theta$ (\textbf{restricted ML}): consistent, asymptotically normal and asymptotically
       efficient \emph{under $H_{0}$}.
 \item $\hat\theta$ (\textbf{unrestricted ML}): consistent and asymptotically normal, \emph{but not
       efficient under $H_{0}$}; it is, however, consistent over a broader range of parameter values
       (whether or not $H_{0}$ holds).
\end{flatlist}

\runin{Note}
\begin{flatnum}
 \item If $H_{0}$ is true, both $\tilde\theta$ and $\hat\theta$ are consistent, and $\tilde\theta$ is more
       efficient. If $H_{0}$ is false, $\tilde\theta$ is inconsistent but $\hat\theta$ is still consistent.
 \item Look at $\hat\theta-\tilde\theta$. If $H_{0}$ is true, $\operatorname{plim}(\hat\theta-\tilde\theta)=0$; if $H_{0}$
       is false, $\operatorname{plim}(\hat\theta-\tilde\theta)\neq0$.
\end{flatnum}
So we create a test statistic based on the contrast $\hat\theta-\tilde\theta$.

\topic{The key result}
Under $H_{0}$,
\begin{equation}\tag{1}
 \boxed{\ \operatorname{var}(\hat\theta-\tilde\theta)=\operatorname{var}(\hat\theta)-\operatorname{var}(\tilde\theta)\ }
\end{equation}
In general
\[ \operatorname{var}(\hat\theta-\tilde\theta)=\operatorname{var}(\hat\theta)+\operatorname{var}(\tilde\theta)-2\operatorname{cov}(\hat\theta,\tilde\theta), \]
so for (1) to hold we need
\begin{equation}\tag{2}
 \operatorname{var}(\tilde\theta)=\operatorname{cov}(\hat\theta,\tilde\theta).
\end{equation}

\topic{Proof of (2)}
Consider the linear combination of $\tilde\theta$ and $\hat\theta$
\[ \theta^{*}=(1-\lambda)\tilde\theta+\lambda\hat\theta . \]
It is also consistent under $H_{0}$, with
\[ \operatorname{var}(\theta^{*})=(1-\lambda)^{2}\operatorname{var}(\tilde\theta)+\lambda^{2}\operatorname{var}(\hat\theta)
   +2\lambda(1-\lambda)\operatorname{cov}(\tilde\theta,\hat\theta). \]
Choose $\lambda$ to minimise $\operatorname{var}(\theta^{*})$:
\[ \frac{\partial\operatorname{var}(\theta^{*})}{\partial\lambda}
   =-2(1-\lambda)\operatorname{var}(\tilde\theta)+2\lambda\operatorname{var}(\hat\theta)+(2-4\lambda)\operatorname{cov}(\tilde\theta,\hat\theta)=0, \]
\begin{equation}\tag{3}
 \Longrightarrow\quad
 \lambda\bigl[\operatorname{var}(\tilde\theta)+\operatorname{var}(\hat\theta)-2\operatorname{cov}(\hat\theta,\tilde\theta)\bigr]
 =\operatorname{var}(\tilde\theta)-\operatorname{cov}(\hat\theta,\tilde\theta).
\end{equation}
By assumption $\tilde\theta$ is asymptotically efficient. Therefore $\lambda=0$, i.e.\ $\theta^{*}=\tilde\theta$,
gives the most efficient estimator in this family, so $\lambda=0$ must solve (3). Then
\[ \operatorname{var}(\tilde\theta)-\operatorname{cov}(\hat\theta,\tilde\theta)=0
   \quad\Longrightarrow\quad\operatorname{var}(\tilde\theta)=\operatorname{cov}(\hat\theta,\tilde\theta). \]
(For vector $\theta$, take $\theta^{*}=\tilde\theta+\lambda A(\hat\theta-\tilde\theta)$ for any fixed $k\times k$ matrix $A$
and any $c$. Efficiency of $\tilde\theta$ means $\operatorname{var}(c'\theta^{*})$ is minimised at $\lambda=0$, so its
derivative there, $2\,c'\operatorname{cov}(\tilde\theta,\hat\theta-\tilde\theta)A'c$, is zero for all $c$ and $A$. Hence
$\operatorname{cov}(\tilde\theta,\hat\theta-\tilde\theta)=0$, i.e.\ $\operatorname{cov}(\hat\theta,\tilde\theta)=\operatorname{var}(\tilde\theta)$.) Substituting (2) into the
general expression gives (1). \hfill$\square$

\topic{The Hausman statistic}
Now look at
\begin{equation}\tag{4}
 (\hat\theta-\tilde\theta)'\bigl[\operatorname{var}(\hat\theta)-\operatorname{var}(\tilde\theta)\bigr]^{-}(\hat\theta-\tilde\theta)
 \ \overset{a}{\sim}\ \chi^{2}_{r},
\end{equation}
where ${}^{-}$ denotes a generalised inverse (see below for why an ordinary inverse does not exist).
Recall
\[ \operatorname{var}(\hat\theta)=\frac1n\,\mathcal I^{-1}(\hat\theta), \]
and, from \emph{MLE~IV},
\begin{align*}
 \operatorname{var}(\tilde\theta)
 &=\frac1n\,\mathcal I^{-1}(\tilde\theta)\Bigl[I-R'(\tilde\theta)\bigl[R(\tilde\theta)\mathcal I^{-1}(\tilde\theta)R'(\tilde\theta)\bigr]^{-1}R(\tilde\theta)\mathcal I^{-1}(\tilde\theta)\Bigr]
 \qquad(I=\text{identity matrix})\\
 &=\frac1n\Bigl[\mathcal I^{-1}(\tilde\theta)-\mathcal I^{-1}(\tilde\theta)R'(\tilde\theta)\bigl[R(\tilde\theta)\mathcal I^{-1}(\tilde\theta)R'(\tilde\theta)\bigr]^{-1}R(\tilde\theta)\mathcal I^{-1}(\tilde\theta)\Bigr].
\end{align*}
Under $H_{0}$ both information matrices estimate the same $\mathcal I(\theta_{0})$, so (evaluating at $\tilde\theta$)
\begin{equation}\tag{5}
 \operatorname{var}(\hat\theta)-\operatorname{var}(\tilde\theta)
 =\frac1n\,\mathcal I^{-1}(\tilde\theta)R'(\tilde\theta)\bigl[R(\tilde\theta)\mathcal I^{-1}(\tilde\theta)R'(\tilde\theta)\bigr]^{-1}R(\tilde\theta)\mathcal I^{-1}(\tilde\theta),
\end{equation}
and the Hausman statistic is
\[ n\,(\hat\theta-\tilde\theta)'\Bigl[\mathcal I^{-1}(\tilde\theta)R'(\tilde\theta)\bigl[R(\tilde\theta)\mathcal I^{-1}(\tilde\theta)R'(\tilde\theta)\bigr]^{-1}R(\tilde\theta)\mathcal I^{-1}(\tilde\theta)\Bigr]^{-}(\hat\theta-\tilde\theta)
   \ \overset{a}{\sim}\ \chi^{2}_{r}. \]
The $k\times k$ matrix in (5) has rank $r$ only (it is sandwiched around the $r\times r$ matrix
$(R\mathcal I^{-1}R')^{-1}$), so for $r<k$ it is singular: this is why a generalised inverse is used, and why
the degrees of freedom are $r$, not $k$.

\topic{Illustration: the linear model with $R\beta=r$}
Let $y=X\beta+u$, $u\sim N(0,\sigma^{2}I)$, with OLS $b=(X'X)^{-1}X'y$, $\operatorname{var}(b)=\sigma^{2}(X'X)^{-1}$, and
restricted least squares (Lesson~6)
\begin{align*}
 b^{*}&=b+(X'X)^{-1}R'\bigl(R(X'X)^{-1}R'\bigr)^{-1}(r-Rb),\\
 \operatorname{var}(b^{*})&=\sigma^{2}(X'X)^{-1}-\sigma^{2}(X'X)^{-1}R'\bigl(R(X'X)^{-1}R'\bigr)^{-1}R(X'X)^{-1}.
\end{align*}
Hence
\begin{align*}
 b-b^{*}&=(X'X)^{-1}R'\bigl(R(X'X)^{-1}R'\bigr)^{-1}(Rb-r),\\
 \operatorname{var}(b)-\operatorname{var}(b^{*})&=\sigma^{2}(X'X)^{-1}R'\bigl(R(X'X)^{-1}R'\bigr)^{-1}R(X'X)^{-1}
 \;\equiv\;\sigma^{2}P .
\end{align*}
One generalised inverse of $P$ is $X'X$ (check: $P(X'X)P=P$), and since $b-b^{*}$ lies in the column space
of $P$ the quadratic form does not depend on which generalised inverse is used. Therefore
\[ \frac1{\sigma^{2}}(b-b^{*})'P^{-}(b-b^{*})
   =\frac1{\sigma^{2}}(Rb-r)'\bigl(R(X'X)^{-1}R'\bigr)^{-1}(Rb-r)\ \sim\ \chi^{2}_{r}, \]
which is exactly the Wald statistic for $R\beta=r$ (and $r$ times the $F$ statistic once $\sigma^{2}$ is replaced
by $s^{2}$). In this simple case the Hausman contrast reproduces the classical test; its value lies in
cases where $\hat\theta$ and $\tilde\theta$ are any efficient/robust pair of estimators, not only ML.

\section{The LM test by an artificial regression}

\topic{The model}
Let
\[ y_{t}=X_{t}\beta+u_{t},\qquad u_{t}=\rho u_{t-1}+\varepsilon_{t},\qquad \varepsilon_{t}\sim\text{i.i.d. }N(0,\sigma^{2}),\qquad t=1,\dots,n, \]
with $X_{t}$ the $1\times k$ row of regressors. The model can be rewritten as
\[ y_{t}=\rho y_{t-1}+X_{t}\beta-\rho X_{t-1}\beta+\varepsilon_{t}
   \;=\;f_{t}(\beta,\rho)+\varepsilon_{t}\;=\;f_{t}(\alpha)+\varepsilon_{t},\qquad\alpha=(\beta,\rho). \]
Define
\[ F_{ti}(\alpha)=\frac{\partial f_{t}(\alpha)}{\partial\alpha_{i}},\qquad t=1,\dots,n,\quad i=1,\dots,p,\quad p=k+1. \]
Therefore $F$ is an $n\times p$ matrix,
\[ F=\begin{bmatrix}
 \dfrac{\partial f_{1}}{\partial\alpha_{1}}&\dfrac{\partial f_{1}}{\partial\alpha_{2}}&\cdots&\dfrac{\partial f_{1}}{\partial\alpha_{p}}\\[8pt]
 \dfrac{\partial f_{2}}{\partial\alpha_{1}}&\dfrac{\partial f_{2}}{\partial\alpha_{2}}&\cdots&\dfrac{\partial f_{2}}{\partial\alpha_{p}}\\[4pt]
 \vdots&\vdots&\ddots&\vdots\\[2pt]
 \dfrac{\partial f_{n}}{\partial\alpha_{1}}&\dfrac{\partial f_{n}}{\partial\alpha_{2}}&\cdots&\dfrac{\partial f_{n}}{\partial\alpha_{p}}
\end{bmatrix}. \]

\topic{Log-likelihood, score and information}
The log-likelihood is
\[ L(\beta,\rho,\sigma)=L(\alpha,\sigma)=\text{const}-n\log\sigma-\frac1{2\sigma^{2}}\sum_{t=1}^{n}\bigl(y_{t}-f_{t}(\alpha)\bigr)^{2}. \]
The first-order conditions are
\begin{align}
 \frac{\partial L}{\partial\alpha_{i}}&=g_{i}(\alpha,\sigma)=\frac1{\sigma^{2}}\sum_{t=1}^{n}\bigl(y_{t}-f_{t}(\alpha)\bigr)F_{ti}(\alpha)=0,\qquad i=1,\dots,p, \tag{6}\\
 \frac{\partial L}{\partial\sigma}&=g_{\sigma}(\alpha,\sigma)=-\frac n\sigma+\frac1{\sigma^{3}}\sum_{t=1}^{n}\bigl(y_{t}-f_{t}(\alpha)\bigr)^{2}=0. \notag
\end{align}
In vector form $g_{\alpha}(\alpha,\sigma)=\frac1{\sigma^{2}}F'(\alpha)\bigl(y-f(\alpha)\bigr)$. The information matrix, with
$\theta=(\alpha,\sigma)=(\beta,\rho,\sigma)$, is
\[ \mathcal I(\theta)=-\E\Bigl[\frac{\partial^{2}L}{\partial\theta\,\partial\theta'}\Bigr]
   =\begin{bmatrix}\dfrac1{\sigma^{2}}F'F&0\\[6pt]0&\dfrac{2n}{\sigma^{2}}\end{bmatrix},
   \qquad
   \mathcal I^{-1}(\theta)=\begin{bmatrix}\sigma^{2}(F'F)^{-1}&0\\[4pt]0&\dfrac{\sigma^{2}}{2n}\end{bmatrix}
   =\begin{bmatrix}\mathcal I^{-1}(\alpha)&0\\0&\mathcal I^{-1}(\sigma)\end{bmatrix}, \]
where $\sigma^{2}(F'F)^{-1}=\sigma^{2}\bigl(F'(\alpha)F(\alpha)\bigr)^{-1}$ from the definition of $F$. (The
off-diagonal block is zero because $\E[\varepsilon_{t}]=0$; the $\sigma$ entry is
$\E[-\partial^{2}L/\partial\sigma^{2}]=-n/\sigma^{2}+3n\sigma^{2}/\sigma^{4}=2n/\sigma^{2}$.)

\topic{The LM statistic}
Let $H_{0}$ impose $r$ restrictions on $\alpha$, and let $\tilde\theta=(\tilde\alpha,\tilde\sigma)$ be the restricted
ML estimator. Then
\[ LM=g'(\tilde\theta)\,\mathcal I^{-1}(\tilde\theta)\,g(\tilde\theta)=g'(\tilde\alpha)\,\mathcal I^{-1}(\tilde\alpha)\,g(\tilde\alpha), \]
since $\mathcal I$ is block diagonal and $g_{\sigma}(\tilde\theta)=0$ ($\tilde\sigma$ solves its own first-order
condition). (Written with the Hessian, $LM=-\tilde g'\tilde H^{-1}\tilde g$; the minus sign disappears once
$-H$ is replaced by the information matrix.) With $\tilde f=f(\tilde\alpha)$ and $\tilde F=F(\tilde\alpha)$,
\begin{align*}
 LM&=\frac1{\tilde\sigma^{2}}\bigl(y-f(\tilde\alpha)\bigr)'F(\tilde\alpha)\;\tilde\sigma^{2}\bigl(F'(\tilde\alpha)F(\tilde\alpha)\bigr)^{-1}\;\frac1{\tilde\sigma^{2}}F'(\tilde\alpha)\bigl(y-f(\tilde\alpha)\bigr)\\
   &=\frac1{\tilde\sigma^{2}}(y-\tilde f)'\tilde F(\tilde F'\tilde F)^{-1}\tilde F'(y-\tilde f).
\end{align*}
But this is just the explained sum of squares from the artificial regression
\begin{equation}\tag{7}
 \frac{y-\tilde f}{\tilde\sigma}=\tilde F\,b+\text{errors}.
\end{equation}
Recall: if $y=X\beta+u$, then $Xb=X(X'X)^{-1}X'y$ and
\[ \mathrm{ESS}=b'X'Xb=y'X(X'X)^{-1}X'X(X'X)^{-1}X'y=y'X(X'X)^{-1}X'y. \]

\runin{To summarise}
$LM=\mathrm{ESS}$ from (7), and $\mathrm{ESS}\overset{a}{\sim}\chi^{2}_{r}$, $r$ = number of restrictions.

\topic{The $nR^{2}$ form}
Also note the following. The residual sum of squares of (7) is
\[ \mathrm{RSS}=\frac1{\tilde\sigma^{2}}(y-\tilde f)'\bigl[I-\tilde F(\tilde F'\tilde F)^{-1}\tilde F'\bigr](y-\tilde f),
   \qquad\text{but}\quad\tilde\sigma^{2}=\frac{(y-\tilde f)'(y-\tilde f)}{n}\ \text{under }H_{0}. \]
Therefore
\[ \frac1{\tilde\sigma^{2}}(y-\tilde f)'(y-\tilde f)=\frac{(y-\tilde f)'(y-\tilde f)}{(y-\tilde f)'(y-\tilde f)/n}=n, \]
and hence
\[ \mathrm{RSS}=\underbrace{\frac1{\tilde\sigma^{2}}(y-\tilde f)'(y-\tilde f)}_{=\,n}
   -\underbrace{\frac1{\tilde\sigma^{2}}(y-\tilde f)'\tilde F(\tilde F'\tilde F)^{-1}\tilde F'(y-\tilde f)}_{=\,LM}
   =n-LM
   \quad\Longrightarrow\quad LM=n-\mathrm{RSS}. \]
Also the (uncentred) total sum of squares of the regressand is
$\mathrm{TSS}=(y-\tilde f)'(y-\tilde f)/\tilde\sigma^{2}=n$, and thus
\[ R^{2}=\frac{\mathrm{ESS}}{\mathrm{TSS}}=1-\frac{\mathrm{RSS}}{\mathrm{TSS}}=1-\frac{n-LM}{n}=\frac{LM}{n}
   \qquad\Longrightarrow\qquad\boxed{\,LM=nR^{2}\,} \]

\section{Example: testing for AR(1) errors, $H_{0}\colon\rho=0$}

From above,
\[ y_{t}=f_{t}(\alpha)+\varepsilon_{t}=\rho y_{t-1}+X_{t}\beta-\rho X_{t-1}\beta+\varepsilon_{t}. \]
Let $H_{0}\colon\rho=0$ (one restriction). With $f_{t}(\alpha)=\rho y_{t-1}+X_{t}\beta-\rho X_{t-1}\beta$,
\begin{align*}
 F_{t\beta_{i}}(\alpha)&=X_{ti}-\rho X_{t-1,i},\qquad i=1,\dots,k,\\
 F_{t\rho}(\alpha)&=y_{t-1}-X_{t-1}\beta .
\end{align*}
Also $\alpha=(\beta,\rho)$ and $\tilde\alpha=(\tilde\beta,0)$, so $\tilde f_{t}=X_{t}\tilde\beta$ and the two equations above become
\[ \tilde F_{t\beta_{i}}=X_{ti},\qquad \tilde F_{t\rho}=y_{t-1}-X_{t-1}\tilde\beta . \]
The artificial regression (7) now becomes
\begin{equation}\tag{8}
 \frac1{\tilde\sigma}\bigl(y_{t}-X_{t}\tilde\beta\bigr)=X_{t}\,c+\bigl(y_{t-1}-X_{t-1}\tilde\beta\bigr)\,d+\text{errors}.
\end{equation}
From (8) we obtain $nR^{2}\overset{a}{\sim}\chi^{2}_{1}$; also a $t$-statistic on $d$ is a valid test of
$H_{0}\colon\rho=0$.

\runin{Note} $\tilde\beta$ is the OLS estimator under $H_{0}\colon\rho=0$ (with $\rho=0$ the model is the classical
normal linear model, so restricted ML of $\beta$ is OLS).

\topic{To summarise: testing for $\rho=0$}
\begin{flatlist}
 \item[Step 1] Run an OLS regression of $y_{t}$ on $X_{t}$ and obtain the residuals $e_{t}$.
 \item[Step 2] Regress $e_{t}$ on $X_{t}$ and $e_{t-1}$, and obtain $nR^{2}\overset{a}{\sim}\chi^{2}_{1}$.
\end{flatlist}
\runin{Note 1} If we do not divide $e_{t}$ by $\tilde\sigma$ we lose nothing, since $R^{2}$ is invariant to rescaling the
regressand.

\runin{Note 2} The regressors in (8) are $X_{t}$ themselves, not the fitted values $X_{t}\tilde\beta$ (these are
the derivatives $\tilde F_{t\beta}$). Because $e$ is orthogonal to $X$, if $X$ contains a constant the
residuals have mean zero and the centred and uncentred $R^{2}$ coincide. For $t=1$ the lag $e_{0}$ is
unavailable: either set $e_{0}=0$ or drop the first observation (asymptotically irrelevant).

\runin{Exercise} Generalise the above procedure to test $H_{0}\colon\rho_{1}=\rho_{2}=\rho_{3}=\cdots=\rho_{p}=0$ if
\[ y_{t}=X_{t}\beta+u_{t},\qquad u_{t}=\rho_{1}u_{t-1}+\rho_{2}u_{t-2}+\cdots+\rho_{p}u_{t-p}+\varepsilon_{t}. \]
\emph{Sketch (added).} Now $f_{t}(\alpha)=X_{t}\beta+\sum_{j=1}^{p}\rho_{j}(y_{t-j}-X_{t-j}\beta)$ and at
$\tilde\alpha=(\tilde\beta,0,\dots,0)$ the derivatives are $\tilde F_{t\beta}=X_{t}$,
$\tilde F_{t\rho_{j}}=y_{t-j}-X_{t-j}\tilde\beta=e_{t-j}$. So regress $e_{t}$ on $X_{t},e_{t-1},\dots,e_{t-p}$;
then $nR^{2}\overset{a}{\sim}\chi^{2}_{p}$ (or use the $F$ test of $d_{1}=\cdots=d_{p}=0$).

\runin{Difference from the GNR?} The regression in (8) has the same regressors as the Gauss--Newton
regression test for AR(1) errors in Lesson~13: $\tilde u_{t}$ on $X_{t}$ and $\tilde u_{t-1}$. The difference is
in how it is motivated and read. The GNR comes from NLS and is read with the usual $t$/$F$ statistic
(which re-estimates $\sigma^{2}$ from the GNR residuals); the artificial regression here comes from the ML
score and information matrix, its regressand is scaled by $\tilde\sigma$, and the statistic is the LM form
$\mathrm{ESS}=nR^{2}$, which uses the \emph{restricted} $\tilde\sigma^{2}$. Under $H_{0}$ they are asymptotically
equivalent.

\clearpage
\begin{transcriptnotes}
 \item Source: \emph{MLE5.pdf}, six photographed pages: pages 1--3 handwritten, headed ``1\,(V)'', ``2'', ``3''
       (Hausman specification testing); pages 4--6 printed sheets numbered 53--55 with handwritten heading
       ``By an artificial Regression: Example of LM test''. Section headings are \textbf{added}. Equation
       numbers: the handwritten part has (1), (2); here the proof's equation (2) is (3) and the condition
       $\operatorname{var}(\tilde\theta)=\operatorname{cov}(\hat\theta,\tilde\theta)$ is labelled (2). The printed sheets number only (8),
       which is kept; (4)--(7) are \textbf{added}. Transposes ($^{T}$) are primes.
 \item Page 1: title ``Supplementary Lecture on Maximum Likelihood. Hausman specification testing (Different
       from W, LM, LR)''. ``consistent over a broader range of estimate'' is read as ``of parameter values''.
       ``Create a test statistic based on $(\tilde\theta-\tilde\theta)$'' (both tildes) is read as $\hat\theta-\tilde\theta$.
 \item Page 2: \textbf{Corrected:} the derivative of $\operatorname{var}(\theta^{*})$ is written with $+2(1-\lambda)\operatorname{var}(\tilde\theta)$;
       differentiating $(1-\lambda)^{2}$ gives $-2(1-\lambda)$. With the corrected sign the handout's next line
       (3) follows; with the handout's sign it would not. The reason $\lambda=0$ must solve (3) and the
       extension to vector $\theta$ are \textbf{added}. In the statistic (4) the handout writes an ordinary inverse;
       the matrix is singular of rank $r$, so a generalised inverse is used (\textbf{corrected}).
 \item Page 3: \textbf{Corrected:} inside the restricted variance the handout writes
       $[R(\tilde\theta)\mathcal I(\tilde\theta)R'(\tilde\theta)]^{-1}$; it must be $[R\,\mathcal I^{-1}R']^{-1}$, as in \emph{MLE~IV}
       (11). The lower half of the page is faint pencil working on the linear model ($b$, $b^{*}$, their
       variances and the contrast $b-b^{*}$); it is partly illegible and is written out cleanly here as
       ``Illustration: the linear model'', with the generalised-inverse step and the identification with the
       Wald statistic \textbf{added}.
 \item Page 4 (sheet 53): the log-likelihood is $L(\beta,\rho,\sigma^{2})=L(\alpha,\sigma^{2})$ in the handout but is written
       and differentiated in $\sigma$; it is labelled $L(\alpha,\sigma)$ here. The handout writes
       $\partial^{2}L/\partial\theta\partial\theta'=-\mathbb H(\theta)$ for the block matrix; that matrix is the
       information matrix $-\E[\partial^{2}L/\partial\theta\partial\theta']$ (\textbf{corrected}, written $\mathcal I$).
       \textbf{Corrected:} the lower-right block of $\mathcal I^{-1}$ is printed $2n/\sigma^{2}$; the inverse is $\sigma^{2}/(2n)$.
       ``Let $H_{0}\colon\alpha=\alpha_{0}$'' is generalised to $r$ restrictions on $\alpha$, as the summary (``$r$ restrictions'')
       and the example ($\rho=0$ only) require.
 \item Page 5 (sheet 54): the printed ``$LM=-g'(\tilde\alpha)H^{-1}(\tilde\alpha)g(\tilde\alpha)=$ ???minus???'' is resolved: with the
       Hessian the minus is needed, with the information matrix (as used in the next line) it is not.
       The top of the sheet is curled; the line $g=\frac1{\sigma^{2}}(y-f(\alpha))'F(\alpha)$ is read from the visible part.
       The struck-through $X'X(X'X)^{-1}$ in the ESS ``recall'' is the cancellation, written out here.
       ``where $r$ -- restrictions (dimension of $\alpha$)'' is read as the number of restrictions (\textbf{corrected}).
       $R^{2}$ here is the uncentred $R^{2}$ (\textbf{added}).
 \item Page 6 (sheet 55): in (8) the first regressor is printed ``$(X_{t}\tilde\beta)c$'' with $\tilde\beta$ crossed out by hand,
       and the printed ``Note 2: Instead of fitted values $X_{t}\tilde\beta$ as a regressor one can simply have the
       $X_{t}$'s'' is struck through; the regressor is $X_{t}$, which is what the derivative gives, and the point is
       kept as Note 2. The $e_{0}$ remark and the exercise sketch are \textbf{added}. The handwritten question
       ``Difference from GNR?'' is answered (\textbf{added}).
 \item Original page map: p.\,1 Hausman setup, result (1); p.\,2 proof of $\operatorname{var}(\tilde\theta)=\operatorname{cov}(\hat\theta,\tilde\theta)$, the
       statistic; p.\,3 $\operatorname{var}(\tilde\theta)$, the contrast's variance, linear-model working; p.\,4 (53) AR(1) model as
       $f_{t}(\alpha)$, $F$, log-likelihood, information; p.\,5 (54) LM as ESS, $LM=nR^{2}$; p.\,6 (55) test for $\rho=0$,
       exercise.
\end{transcriptnotes}

\end{document}
